\chapter{LITERATURE REVIEW}
\hspace{0.5cm}
Maintaining accurate and up-to-date software documentation remains one of the greatest challenges in contemporary software development. As software systems continue to evolve through continuous integration and agile methodologies, documentation created manually often becomes inconsistent with the source code. This mismatch between code and its documentation reduces code legibility and makes it challenging for new developers joining a codebase. To address these problems, researchers have recently made efforts toward automating document generation using artificial intelligence. In particular, combining structural program analysis with LLM has emerged as a promising solution. This chapter reviews the most relevant literature, organized around key themes including code representation strategies, repository-level documentation approaches, empirical evaluations of language models, and practical deployment considerations.

A fundamental challenge in automated code documentation is how to represent source code in a way that captures its meaning and structure for machine processing. Early approaches treated code as simple sequences of tokens, but this ignored the rich syntactic structure that distinguishes programming languages from natural language. Researchers have increasingly recognized that AST offer a more powerful representation, as they break code down into meaningful components such as functions, classes, and control statements while preserving hierarchical relationships. However, simply serializing ASTs into sequences for neural network processing loses important positional relationships between nodes, including hierarchical parent-child relations and sibling order. To address this limitation, researchers have proposed tree positional embeddings that inject both vertical and horizontal positional information for each node in the AST, along with specialized syntax-tree attention mechanisms that focus on syntactically important code tokens. This enhanced AST encoding has demonstrated consistent improvements across multiple programming languages and evaluation metrics, outperforming state-of-the-art models by significant margins on both Java and Python benchmarks \cite{Song2024}. Beyond traditional AST representations, researchers have explored treating code as graphs to capture relationships that tree structures cannot adequately represent. When ASTs are transformed into graph representations where nodes correspond to syntactic units and edges represent hierarchical relationships, graph neural networks can extract structural features through neural message passing with neighborhood aggregation. This graph-based perspective preserves connectivity information that tree representations miss, such as complex dependency relationships and recursive patterns, and has achieved high accuracy in code understanding tasks \cite{Rozi2023}. Additionally, incorporating code metrics such as complexity measures, API (Application Programming Interface) popularity, and Halstead metrics into documentation models has been shown to improve generation quality, suggesting that structural characteristics beyond raw syntax contribute meaningfully to code comprehension \cite{Mostafavi2024}.

The integration of multiple complementary code representations has emerged as another important direction for improving documentation quality. Rather than relying on a single representation, combining token sequences with type-augmented ASTs and control flow graphs provides a more complete picture of code semantics. Type augmentation enriches ASTs with type information about variables and expressions, while control flow graphs capture program execution order and branching structure. Hierarchical attention mechanisms have proven effective for integrating these diverse representations, with one attention layer operating at the token level within statements and another layer operating at the statement level within functions. This hierarchical structure mirrors the natural hierarchy of code, where tokens form statements and statements form functions. Reinforcement learning has also been employed to address the exposure bias problem in sequence generation, where models trained to predict the next word given ground truth context must rely on their own potentially incorrect predictions during actual use. Actor-critic networks jointly optimize word selection and evaluate overall sequence quality, leading to substantial improvements over maximum-likelihood training approaches \cite{Wang2020}.

Documenting entire software repositories presents challenges beyond those of individual function summarization. Understanding a full software project requires comprehending inter-file dependencies, system architecture, and how different modules collaborate to implement overall functionality. One framework addresses this by constructing a project tree that reflects the repository's directory structure, then extracting reference relationships between code objects to form a directed graph that captures semantic call relationships. These reference relationships serve as natural contextual learning examples, teaching language models how different parts of the codebase interact. The framework generates structured documentation covering functionality, parameters, code descriptions, considerations, and examples, using prompt templates that incorporate project trees, code snippets, reference relationships, and metadata \cite{CodeDocAgent2025}. A feature-oriented approach to repository summarization offers an alternative to directory-based organization. Rather than relying solely on file structure, this method clusters methods and files by both dependency relationships and semantic similarity to identify functional features that naturally span multiple files. The clustering process uses community detection algorithms on graphs where nodes represent code elements and edges represent both structural dependencies and semantic similarities. Once features are identified, large language models generate descriptive documentation for each feature, and traceability links are preserved between features and the code elements that implement them. This feature-oriented organization has been shown to increase the rate of completely covered features in manual documentation and significantly improve the recall of traceability links between documentation and code \cite{Zhu2025}.

Beyond generating documentation directly from code, researchers have also explored the related task of generating code examples to complement existing documentation. This work recognizes that good documentation includes not only descriptions of what code does but also concrete examples of how to use it. Using language models pre-trained on both natural and programming languages, researchers have generated code examples from source code and documentation inputs. The effectiveness of this approach has been demonstrated through metrics such as passability, measuring whether generated examples execute without error, and relevance, measuring whether examples properly address the target method and documentation. A particularly significant finding is that incorporating compiler error logs into the input improves passability substantially, as the error messages provide feedback that allows the model to correct mistakes in subsequent attempts. This error-feedback mechanism demonstrates the value of iterative refinement based on execution results \cite{KhanUddin2023}.

Empirical evaluation of different language models for documentation tasks has provided valuable guidance for model selection. Comparative studies have examined models such as GPT-2 (Generative Pre-trained Transformer–version 2) and RoBERTa (Robustly Optimized BERT Pretraining Approach ), finding that RoBERTa consistently outperforms GPT-2 in both accuracy and loss metrics. The superior performance of RoBERTa is attributed to its ability to detect semantic connections in code, achieving near-perfect accuracy on documentation generation tasks while maintaining dramatically lower training and validation losses. These results highlight the importance of model architecture choices and the value of fine-tuning with domain-specific datasets for documentation generation \cite{Thota2024}. Similar comparative work on T5 (Text-to-Text Transfer Transformer) and BART (Bidirectional and Auto-Regressive Transformer) has shown that it generates more concise and context-aware comments, making it more reliable for capturing the single intent of code snippets. However, T5 offers advantages in training efficiency, presenting a trade-off between output quality and computational resources. The use of multiple evaluation metrics including BLEU (Bilingual Evaluation Understudy), ROUGE (Recall-Oriented Understudy for Gisting Evaluation), METEOR (Metric for Evaluation of Translation with Explicit ORdering), and Smoothed BLEU provides a comprehensive framework for assessing documentation quality across different dimensions \cite{GhaleDabbagh2025}.

The practical deployment of automated documentation systems in real-world development workflows has also received attention. One system integrates AST generation from a compiler with a large language model within a continuous integration and continuous delivery pipeline, triggered on every code push to automatically generate updated documentation. This integration with version control systems ensures that documentation remains synchronized with code changes, addressing the documentation debt problem directly. A survey of experienced developers found that concise documentation was strongly preferred over verbose output, with developers rating terseness as more important for comprehensibility than additional detail. This finding has important implications for prompt design and model configuration in automated documentation systems, suggesting that generating shorter, focused documentation may be more valuable than producing comprehensive but lengthy explanations \cite{Procko2023}.

In summary, the reviewed literature establishes that effective automated code documentation requires careful attention to code representation, with AST-based, graph-based, and metric-enriched approaches each contributing valuable perspectives. Repository-level documentation benefits from both directory-based and feature-oriented organization strategies. Empirical studies provide guidance on model selection, with different models offering trade-offs between output quality and computational efficiency. Practical deployment in CI/CD (Continous Integration/ Continous Deployment) pipelines has demonstrated feasibility, with user preferences favoring concise documentation. The proposed AI-Powered Code Document Generator builds on these findings by combining Tree-sitter-based AST parsing across fourteen programming languages, NetworkX-based graph analysis for capturing function call relationships, and a locally-hosted LLM to generate comprehensive documentation including README files, architectural descriptions, file-wise explanations, and recent change histories without modifying source code.
\par Even though all the above-mentioned methods are required to create the proposed system, their implementation needs a certain environment with resources and tools prepared beforehand. That is why the proper setting up of the development environment will be important and this point is described in the next chapter.